% !TeX root = ../main.tex
\section{Implementation}
\label{sec:implementation}

Every theorem of the previous sections is mechanized in the Lean programming language~\cite{lean}, using Mathlib's formalization of measure theory and probability~\cite{mathlib}, in particular the Giry monad, Markov kernels and their disintegration, as well as the standard continuous and discrete distributions. Most of the development, including the definitions and theorem statements, was written by large language models (Claude Fable 5.1, Claude Opus 5.5 and GPT 6 Astra). Since Lean's kernel checks the proofs but not whether a statement says what this paper says, we separated the code into two parts:

\begin{enumerate}
\item \textbf{Specification}: the definitions in \texttt{Spec/}, such as the syntax, typing rules, semantics and determinization, and the full statement of every theorem in \texttt{Theorems.lean}, about 1,600 lines in total. We reviewed both line by line against \cref{sec:language,sec:type-system,sec:soundness}, and reviewed every later change to them.
\item \textbf{Proofs and procedures}: about 26,000 lines in \texttt{Proof/} that prove the statements, and the executable code that some statements are about, such as type inference. \texttt{Theorems.lean} proves each theorem by a single reference to this part. Its last command fails the build if a theorem uses an axiom other than \texttt{propext}, \texttt{Classical.choice} and \texttt{Quot.sound}, or if its statement depends, through anything other than proofs, on a declaration in \texttt{Proof/}.
\end{enumerate}

To trust the tool's results, one has to rely on Lean's kernel, these three axioms, the specification and the Mathlib definitions it uses, and the Lean compiler for the output of type inference. Exact results come with a certificate that the kernel checks against the specification. The parser, the elaborator and the floating-point interpreter are unverified.

The complete compilation, verification and execution pipeline is outlined in \cref{fig:architecture}.

% !TeX root = ../main.tex
\begin{figure}[htb]
\centering
\begin{tikzpicture}[
  x=\textwidth, y=1cm,
  font=\small,
  stage/.style={draw=black!45, fill=black!3, rounded corners=2pt,
    text width=.18\textwidth, minimum height=12mm, align=center, inner sep=4pt},
  verified/.style={stage, draw=teal!65!black, fill=teal!7, line width=.8pt},
  specification/.style={stage, draw=orange!70!black, fill=orange!8, line width=.8pt},
  compilation/.style={minimum height=14mm},
  data/.style={text width=.18\textwidth, align=center, inner sep=3pt},
  flow/.style={-{Stealth[length=4pt,width=3pt]}, draw=black!65, line width=.6pt},
  note/.style={font=\footnotesize, text=black!70, align=center},
  heading/.style={font=\small\bfseries, anchor=west}
]
\node[heading] at (0,2.85) {Compilation};
\draw[black!20] (0,2.55) -- (1,2.55);
\node[data, text width=.09\textwidth] (file) at (.053,.45) {\texttt{.det}\\program};
\node[stage, compilation, text width=.15\textwidth] (front) at (.23,.45) {Parse and\\elaborate\leanref{frontend}};
\node[verified, compilation, text width=.15\textwidth] (infer) at (.44,.45) {Infer types\\and E/G\leanref{inference}};
\node[specification, compilation, text width=.155\textwidth] (det) at (.672,-1.05) {Determinize\leanref{determinization}};
\node[data, text width=.14\textwidth] (source) at (.88,1.95) {Annotated\\source $e_{\mathrm{src}}$};
\node[verified, compilation, text width=.14\textwidth] (soundness) at (.88,.45)
  {Soundness theorems\leanref{expectation-preservation}};
\node[data, text width=.14\textwidth] (target) at (.88,-1.05) {Transformed\\program $e_{\mathrm{det}}$};
\draw[flow] (file) -- (front);
\draw[flow] (front) -- (infer);
\draw[draw=black!65, line width=.6pt] (infer.east) -- (.555,.45);
\draw[flow] (.555,.45) |- (det.west);
\draw[flow] (.555,.45) |- (source.west);
\draw[flow] (det) -- (target);
\draw[flow] (source) -- (soundness);
\draw[flow] (target) -- (soundness);

\node[heading] at (0,-2.35) {Program evaluation};
\draw[black!20] (0,-2.65) -- (1,-2.65);
\node[data] (program) at (.10,-5.15) {Program $e$};
\node[note] at (.10,-5.60) {$e\in\{e_{\mathrm{src}},e_{\mathrm{det}}\}$};
\node[specification] (semantics) at (.365,-3.50) {Measure-theoretic\\semantics\leanref{output-measures}};
\node[data, text width=.16\textwidth] (measure) at (.88,-3.50) {Output measure\\$\law e$};
\draw[flow] (semantics) -- (measure);
\draw[flow] (measure.east) -- (.99,-3.50) |- (soundness.east);

\node[stage] (runtime) at (.365,-5.15) {Floating-point\\interpreter\leanref{interpreter}};
\node[data, text width=.16\textwidth] (estimates) at (.88,-5.15) {Empirical mean\\and variance};
\draw[flow] (runtime) -- (estimates);

\node[verified] (explore) at (.365,-6.80) {Finite graph\\construction\leanref{graph-construction}};
\node[verified] (solve) at (.635,-6.80) {Exact rational\\solver\leanref{rational-solver}};
\node[data, text width=.16\textwidth] (answer) at (.88,-6.80)
  {Exact mean\\and variance};
\draw[draw=black!65, line width=.6pt] (program.east) -- (.235,-5.15);
\draw[draw=black!65, line width=.6pt] (.235,-3.50) -- (.235,-6.80);
\draw[flow] (.235,-3.50) -- (semantics.west);
\draw[flow] (.235,-5.15) -- (runtime.west);
\draw[flow] (.235,-6.80) -- (explore.west);
\draw[flow] (explore) -- (solve);
\draw[flow] (solve) -- (answer);

\node[stage] (storm) at (.365,-8.50) {Storm solver};
\node[verified] (check) at (.635,-8.50) {Certificate\\checker\leanref{certificate-checker}};
\draw[flow] (explore.south) -- (storm.north);
\draw[flow] (storm) -- (check);
\draw[flow,shorten >=5pt] (check.east) -| (answer.south);

\node[inner sep=0pt] at (.50,-9.80) {
  \begin{tikzpicture}
  \node[specification, minimum width=5mm, minimum height=3mm, text width=0pt,
    inner sep=0pt,anchor=west] (specificationSwatch) at (0,0) {};
  \node[note,anchor=west,inner sep=0pt] (specificationLegend) at ([xshift=5pt]specificationSwatch.east) {Specification};
  \node[verified, minimum width=5mm, minimum height=3mm, text width=0pt,
    inner sep=0pt,anchor=west] (verifiedSwatch) at ([xshift=18pt]specificationLegend.east) {};
  \node[note,anchor=west,inner sep=0pt] (verifiedLegend) at ([xshift=5pt]verifiedSwatch.east) {Proved or checked in Lean};
  \node[stage, minimum width=5mm, minimum height=3mm, text width=0pt,
    inner sep=0pt,anchor=west] (unverifiedSwatch) at ([xshift=18pt]verifiedLegend.east) {};
  \node[note,anchor=west,inner sep=0pt] at ([xshift=5pt]unverifiedSwatch.east) {Unverified};
  \end{tikzpicture}
};
\end{tikzpicture}
\caption{Implementation architecture. Box styles say how each stage is justified; Lean marks link to the code.}
\label{fig:architecture}
\Description{Compilation parses and elaborates the program without verification,
then runs verified type and E/G inference, which branches up to the annotated
source and down through determinization to the transformed program. Between
these two programs sit the soundness theorems.
One program branches to formal measure semantics, floating-point sampling,
and finite graph construction, which is checked against the semantics. The measure semantics yields an output
measure, which feeds back up into the soundness theorems: they relate the output
measures of the annotated source and the transformed program. The graph feeds either a
verified exact rational solver or Storm's exact model checking of expected
accumulated rewards, followed by a verified certificate checker.
Both routes produce exact mean and variance conditioned on returning, when return probability is positive.}
\end{figure}


\paragraph{Parse and elaborate} First, the source program is parsed and desugared. This part is unverified.

\paragraph{Infer types and $\E/\G$} Type inference then fills in every omitted $\E/\G$ annotation (the Lean development calls $\E$ and $\G$ affinities). By \cref{thm:inference}, it fails only if no completion exists and otherwise returns the greatest completion, which has an $\E$ at every site where some completion has one. The result is a typed program $e_{\mathrm{src}}$.

\paragraph{Determinize} The determinization function of \cref{def:determinization} walks through the syntax tree and replaces every $\E$-mode sample by a mean site, marked $\M$. The pipeline runs the same Lean definition that the soundness theorems are stated about, with rational instead of real literals. The output is a well-typed, determinized program $e_{\mathrm{det}}$.

\paragraph{Soundness theorems} The theorems of \cref{sec:soundness} relate the output laws of $e_{\mathrm{src}}$ and $e_{\mathrm{det}}$ under the hypotheses stated there.

\paragraph{Measure-theoretic semantics} These output laws are given by the semantics of \cref{fig:semantics}, which we define in Lean with the Giry monad. They are measures on the reals, which Lean cannot compute, so this semantics is not executed.

\paragraph{Floating-point interpreter} To sample, we also provide an interpreter that evaluates programs with floating-point numbers and pseudo-random samplers. This part is unverified: we have no formal model that relates floating-point arithmetic and the samplers to the real-valued semantics.

\paragraph{Finite graph construction} The tool explores the machine states of the program, each an expression, an environment and a continuation with rational numbers, and merges equal states. Exploration succeeds if finitely many states are reachable and the only random draws left are Bernoulli and discrete ones. With \texttt{-{}-additive}, numbers added to the result of a pending call become rewards on transitions, which keeps loops that accumulate a sum finite. The return probability and the first and second moments then solve linear equations over the resulting Markov chain.

\paragraph{Exact rational solver} A solver written and verified in Lean solves the equations with rational arithmetic.

\paragraph{Storm solver} Alternatively, Storm~\cite{hensel2022storm} solves them, called through its Python interface on an exact rational model.

\paragraph{Certificate checker} Both routes write a Lean file with the chain and a certificate that consists of the solution, the states from which the program can no longer return, and for every other state a path to a terminal or non-returning state, which makes the solution unique. For Storm, our wrapper adds the states and paths. Lean's kernel checks the file, which proves that the explored program is domain-safe, that its output law is that of the chain, and that the certified numbers are its exact return probability, first and second moments, and rejection probability.

\paragraph{Exact mean and variance} Both routes compute the exact probability that the explored program returns a result, and the mean and variance of the result conditioned on that event. If the explored program is $e_{\mathrm{det}}$ and $e_{\mathrm{src}}$ is domain-safe with a defined expectation, \cref{thm:expectation} guarantees that $e_{\mathrm{src}}$ has the same mean.
